\chapter*{Πρόλογος}
\addcontentsline{toc}{chapter}{Πρόλογος}

Το παρόν σύγγραμμα απευθύνεται σε προπτυχιακούς φοιτητές που παρακολουθούν ένα πρώτο πανεπιστημιακό μάθημα κυβερνοασφάλειας, καθώς και σε επαγγελματίες που επιθυμούν να εδραιώσουν την κατανόησή τους για το πεδίο σε στέρεη εννοιολογική βάση. Προϋποθέτει βασική εξοικείωση με υπολογιστές και δίκτυα, όχι όμως προηγούμενη μελέτη της ασφάλειας.

Το βιβλίο οργανώνεται σε δεκατρία κεφάλαια, ομαδοποιημένα σε τέσσερα μέρη. Η σειρά είναι συνειδητή. Ξεκινά από τις έννοιες και τις αρχές, περνά στο λειτουργικό σύστημα ως το πρώτο απτό όριο ασφάλειας και στη συνέχεια μελετά τον αντίπαλο: ποιος επιτίθεται, με ποια εργαλεία και μέσα από ποια κανάλια. Μόνο αφού ο αναγνώστης κατανοήσει τι πρέπει να προστατευθεί, το βιβλίο εισάγει τους μηχανισμούς προστασίας —την κρυπτογραφία, τη διαχείριση ταυτότητας, την ασφαλή μηχανική λογισμικού και τον έλεγχο ασφάλειας. Το τελευταίο μέρος στρέφεται στις λειτουργίες: την ανίχνευση και την απόκριση σε περιστατικά, την ανάλυση κακόβουλου κώδικα και αποδεικτικών στοιχείων και τη διακυβέρνηση της ασφάλειας ως οργανωσιακής λειτουργίας.

Κάθε κεφάλαιο ξεκινά με σύντομη εισαγωγή και με τα μαθησιακά αποτελέσματα, αναπτύσσει το αντικείμενό του σε αριθμημένες ενότητες και ολοκληρώνεται με γλωσσάριο βασικών όρων, σύνοψη και ερωτήσεις ανασκόπησης. Οι πίνακες συμπυκνώνουν συγκρίσεις που γίνονται ευκολότερα κατανοητές όταν παρουσιάζονται παράλληλα, ενώ τα πλαίσια με παραδείγματα συνδέουν τις αφηρημένες έννοιες με ρεαλιστικές καταστάσεις. Οι παραπομπές μεταξύ κεφαλαίων είναι συχνές, επειδή η ασφάλεια είναι ένα διασυνδεδεμένο πεδίο: μια αδυναμία σε ένα επίπεδο συνήθως αξιοποιείται μέσω κάποιου άλλου.

Το βιβλίο εκδίδεται παράλληλα στα αγγλικά και στα ελληνικά. Το ελληνικό κείμενο δεν αποτελεί κατά λέξη μετάφραση· είναι επιμελημένη απόδοση που ακολουθεί τις συμβάσεις του ελληνικού ακαδημαϊκού λόγου και χρησιμοποιεί συνεπή, καθιερωμένη ορολογία. Όπου δεν υπάρχει ευρέως αποδεκτός ελληνικός όρος, ο αγγλικός όρος δίνεται σε παρένθεση κατά την πρώτη εμφάνισή του.
